\section{Implementing the two exact binary64 profiles}\label{app:implementation-guide}

A reference checker can turn every materialized binary64 operand $f$ into a Python \id{Fraction.from_float(f)} and perform each operation on reduced fractions \autocite{pythonFractions2026}. This also specifies a direct test of a faster backend. Given the same operands and routes, compare feasibility, the objective's reduced numerator and positive denominator, and, for TD routes, both the minimum duration and the earliest departure attaining it. A displayed decimal cannot distinguish all these outcomes. The two profiles permit different shortcuts because their route computations use different arithmetic operations.

\subsection{Static routes on a common integer scale}\label{app:static-integers}

\paragraph{Decode each operand once.}
The static checker sums arc costs, propagates times by addition and maximum, and tests inclusive bounds. Multiplying every arc cost, service duration and time-window bound by the same positive power of two preserves all these operations and comparisons. We can choose that scale exactly from the bits of the input. For a nonzero, nonnegative binary64 value, let $E$ and $F$ denote its stored exponent and fraction fields. Recover an integer $M$ and an exponent $e$ such that $f=M2^e$ as follows.
\[
 (M,e)=
 \begin{cases}
 (2^{52}+F,\ E-1023-52), & 1\leq E\leq 2046,\\
 (F,\ -1074), & E=0.
 \end{cases}
\]
The second line covers subnormals. Both signed zeros give zero units, irrespective of the chosen scale. The checker rejects a nonfinite or strictly negative cost or temporal operand. For positive $M$, strip its $k$ trailing zero bits and set $m=M/2^k$ and $h=e+k$. Among all nonzero operands, choose $q$ as the largest of $0$ and the values $-h$ (use $q=0$ if all operands vanish). The integer representing $f$ in units of $2^{-q}$ is exactly $m\,2^{h+q}$. Since $h+q\geq0$, an integer left shift suffices even when a floating-point scale $2^q$ would overflow. The C++ checker implements this decomposition directly.

\paragraph{Evaluate a route in integer units.}
Consider 2 customers, a depot opening at 0, and successive arc costs $3/8$, $5/4$ and $1/2$. The first customer opens at $1/2$ and requires $1/4$ unit of service. The second closes at 2 and requires zero service. Take all remaining temporal bounds to be integral and nonbinding and all remaining arc costs to be integral. The common scale is $2^3$. The vehicle reaches the first customer at tick 3, waits until tick 4 and leaves at tick 6. It reaches the second customer at tick $6+10=16$, exactly its closing time. It returns to the depot at tick 20. The static objective counts travel alone, so the checker returns $(3+10+4)/8=17/8$. In general, if the sum in ticks is $U$, reduce $U/2^q$ before writing the result. Structural integers such as capacity and demand keep their original units.

\paragraph{Choose the accumulator width before evaluating submissions.}
A finite binary64 input can require more than 128 bits on the common scale. For example, a matrix containing both $2^{-1074}$ and 1 needs $q=1074$, so its entry for 1 requires 1,075 bits. Even if every entry fits, an objective or arrival time can exceed the largest single operand. Let $n$ be the number of customers, and let $A$, $S$, $E_{\max}$ and $L_{\max}$ be the largest arc cost, service duration, opening time and closing time in ticks. Use zeros for absent temporal fields. The prototype computes the bound in arbitrary-precision integers as
\[
 B=\max\{2nA,\ E_{\max}+(n+1)(A+S),\ L_{\max}\}.
\]
A solution with $k$ nonempty routes has $n+k\leq2n$ arcs. A route has at most $n+1$ arcs, and waiting can move its clock forward to an opening time of at most $E_{\max}$. Starting from its last wait, or from the depot opening if there was none, leaves at most $n+1$ arc and service increments. The time bound therefore covers every intermediate sum in a structurally valid route, including one that will miss a deadline.

For $B\leq2^{127}-1$, the prototype uses unsigned 128-bit storage and reserves the high bit. It checks each shift before executing it and checks $b\leq (2^{127}-1)-a$ before adding $a+b$. A wider bound selects Boost arbitrary-precision integers at preparation. A checked addition that exceeds the selected range triggers a retry in arbitrary precision. A backend without that fallback must signal an arithmetic range error so the caller can choose another implementation. Treating that error as route infeasibility could discard a valid incumbent. For example, a route may have finite individual arcs and an exact total greater than the largest finite binary64 value.

\paragraph{Pay for preparation where it is reused.}
With a prepared integer matrix, each route leg requires an indexed load and an integer addition. With conversion on access, the checker shifts the binary64 value when the route visits it. Both paths use the same $q$ and return the same reduced result. The measured C++ prototype owns both the binary64 and integer matrices, including when its on-access kernel is selected. Its on-access timings therefore give no estimate of memory savings from omitting the integer matrix. A Python implementation can likewise prepare integer ticks using \id{float.as_integer_ratio()} and Python's arbitrary-precision integers. When all operands are integral, $q=0$ and these ticks are ordinary units. The \id{binary64-exact-v1} identifier still describes their interpretation.

\subsection{TD routes with arbitrary rationals}\label{app:td-rationals}

TD interpolation introduces divisions. If the endpoints of an arrival-time function (ATF) have binary64 coordinates $(0,1)$ and $(6,11)$, its slope is $5/3$. Even an integer departure at time 1 then gives arrival $8/3$. To find which initial departure reaches a window or an arc knot at a particular time, a checker must also invert an affine function. A common power-of-two scale cannot represent all these intermediates. Both C++ implementations convert each finite binary64 operand exactly to a rational using the \href{https://gmplib.org/manual/Rational-Conversions}{GNU Multiple Precision Arithmetic Library (GMP)}. Their output consists of reduced integer pairs for duration and departure.

\paragraph{Compose the route prefix.}
After each visit, retain the vehicle clock $P(t)$ as a function of the initial departure $t$. Each retained knot holds its departure, the attained clock reading and the right-hand limit. If $a<b$ are consecutive knots, $u=P(a^+)$ and $v=P(b)$, then a boundary at $u<z<v$ is reached at
\[
 t_z=a+(z-u)\frac{b-a}{v-u}.
\]
Inserting $t_z$ when $z$ is a knot of the next arc splits the prefix exactly where its formula changes. FIFO puts the encountered knots in increasing clock order, so the algorithm scans forward through the next arc as it processes the prefix. At a jump it evaluates the attained ordinate at the knot. A prefix constant at the jump continues to use that same ordinate, while an increasing prefix that moves immediately past the jump uses the arc's right-hand limit. At a customer, trim the departure domain at the closing bound, insert a cut if the incoming clock crosses the opening time, then apply waiting and service. Before every arc, restrict the preceding clock to its departure horizon. A depot return deadline trims the final prefix once more.

\paragraph{Sweep through initial departures.}
The second algorithm visits the same kinds of boundaries in initial-departure order. Evaluate the route at the left endpoint. From a feasible $t$, trace a small positive displacement $h$ through the route. On the current open interval each clock has the form $u+rh$, where $u$ is its right limit and $r\geq0$ its rate with respect to initial departure. A positive rate reaches the next time boundary $z$ after $(z-u)/r$. Multiply the rate by the next ATF slope at a travel step, and set it to zero when the vehicle waits until a fixed opening time. Record the smallest positive step allowed by an arc knot, a waiting threshold, a horizon or a deadline. The final clock remains affine until this first event. Its value at the event follows by left continuity. A jump can make every departure immediately to its right infeasible even when the event itself is feasible.

Both algorithms consider feasible endpoints, including a singleton domain. On an open affine piece the duration is monotone or constant. The lower value at an upward jump is attained, so no right limit can improve the best attainable value at that departure. Processing cuts in increasing order and replacing the best pair only for a strictly smaller duration retains the first departure on a minimizing plateau. The Python \id{Fraction} reference independently partitions the domain by boundary preimages, evaluates the forward route at the cuts, and checks the one-sided limits. A Python event-sweep variant uses the same rational type with the second traversal.

\paragraph{A jump reached after waiting.}
Take one customer with service window $[2,6]$, zero service and no depot window. Both arcs have departure horizon $[0,6]$. Use the outbound ATF above and give the return arc knots $(0,1)$, $(2,3)$, $(2,4)$ and $(6,8)$. The outbound arrival $1+5t/3$ first reaches 2 at $t=3/5$ and reaches 6 at $t=3$. The feasible initial departures are therefore $[0,3]$. For $0\leq t\leq3/5$, waiting fixes the customer's departure at 2. The attained return-arc ordinate is 3 throughout this plateau, including its right endpoint. For $3/5<t\leq3$, the return arc adds 2 time units. Thus for every feasible departure the duration is
\[
 \delta_{\mathbf r}(t)-t=
 \begin{cases}
 3-t, & 0\leq t\leq3/5,\\
 3+2t/3, & 3/5<t\leq3 .
 \end{cases}
\]
Hence the reported duration is $12/5$ and the earliest departure is $3/5$. Prefix composition inserts the exact preimage $3/5$ of the return-arc knot. The event sweep reaches it from 0 when waiting ends. A traversal that uses the upper ordinate 4 while the vehicle waits at 2 changes the optimum.

\paragraph{Scope of the native implementations.}
The static prototype checks complete solutions. Demands and capacities must be nonnegative integers representable both in signed 64-bit storage and exactly in binary64. Its input adapter enforces zero depot service. The TD prototype checks one prepared route at a time. The calling solution evaluator checks coverage, fleet and capacity, sums the exact durations and adds any vehicle charge. Either native TD routine can therefore use the Python reference's aggregation and serialization.


\paragraph{Comparison on stored routes.}
At \href{https://github.com/ANR-MAMUT/MAMUT-routing/tree/6316c449d5a026c9ac2969a6ce81350dc6f73434}{collection revision \texttt{6316c449}}, both native static kernels matched the Python reference on all 1,976 incumbent costs and 56,620 route travel sums. Every native check used fixed-width arithmetic on these data. Separate boundary checks exercised the arbitrary-precision path. On the same revision the two GMP algorithms and the Python sweep agreed with the Python partitioning reference on all 3,072 TD solution costs, and separately on every minimum duration and earliest departure across 61,992 routes. Exact forward traversal at each reported departure recovered its duration.

On Grid'5000, median static checking times per solution with a prepared matrix were 0.0222~ms for C++ and 1.1134~ms for Python fractions. Including matrix preparation gave 1.0843~ms and 80.1440~ms. Median TD solution times were 29.20~ms for GMP composition, 30.11~ms for the GMP sweep and 5,846.31~ms for Python partitioning. TD times include route preparation, optimization and exact aggregation. File loading and result serialization are excluded in both experiments.
